% titel: Möbelrücken: Kippen unter Schubkraft
% skript: Kap. Statik, Drehmoment, Reibung

Ein Klotz (homogene Masse $m$, Breite $B$, Höhe $H$) wird auf horizontalem
Boden mit einer Kraft $F$ geschoben, die auf der Höhe $h$ angreift. Bei
konstanter Gleitgeschwindigkeit gilt für die Schub- und die Reibungskraft
ein Kräftegleichgewicht; gleichzeitig darf das Drehmoment zum Kippen nicht
überhandnehmen. Gesucht ist die maximale Höhe $h$, bei der der Klotz gerade
noch nicht umkippt.

\begin{werkzeugkasten}[Werkzeuge]
\textbf{Kippbedingung:} Kurz vor dem Kippen dreht der Körper um die vordere
untere Kante (den Kipp-Punkt $P$). Dort greifen dann Normal- und
Reibungskraft an (sie wandern an die Kontaktkante). Um $P$:
\begin{equation}
  \text{Kipp-Drehmoment der Schubkraft}:\; F\cdot h,
  \qquad
  \text{stabilisierendes Gewicht}:\; m g\cdot \frac B2.
\end{equation}

\textbf{Gleitreibung:} bei konstanter Geschwindigkeit gilt in horizontaler
Richtung $F = \mu\, mg$ (die Reibung gleicht die Schubkraft aus).

\textbf{Nicht kippen:} Das kippende Drehmoment darf das stabilisierende nicht
überschreiten:
\begin{equation}
  F\cdot h \le m g\,\frac B2.
\end{equation}

\textbf{Konstanten:} $\mu=\qty{0.34}{}$, $B=\qty{90}{cm}=\qty{0.90}{m}$,
$g=\qty{9.81}{m/s^2}$.
\end{werkzeugkasten}

\begin{werkzeugkasten}[Vorgehen in allgemeiner Form]
$F=\mu mg$ einsetzen, in die Kippbedingung stellen und nach $h_\mathrm{max}$
auflösen; die Höhe $H$ des Klotzes spielt dabei keine Rolle, solange
$h\le H$. Die Frage "wo greifen die Reibungskräfte an?" klären: kurz vor
dem Kippen am Kipp-Punkt $P$.
\end{werkzeugkasten}

\begin{achtung}[Typische Fallen]
\begin{itemize}
  \item Nicht $F$ willkürlich wählen: bei konstanter Geschwindigkeit ist
       $F=\mu mg$.
  \item Der Kipp-Punkt ist die \emph{vordere} untere Kante (Bewegungsrichtung);
       das Gewicht wirkt im Schwerpunkt, Hebelarm $B/2$.
  \item Die Höhe $H$ geht nur über $h\le H$ ein; $h_\mathrm{max}$ hängt von
       $B$ und $\mu$ ab.
\end{itemize}
\end{achtung}